% ============================================================
% f2015tjl_q03_flowchart.tex —— 2015 天津卷（理科）第 3 题 程序框图
% ------------------------------------------------------------
% 坐标口径：图内画布 pt，原点 = 源页基准点，x 右正、y 上正。
%   来源：.work/phase3-b11/figure-inventory.md（真实页 pt，原点左上、y 向下）
%   换算：X_tikz = x_page − 179.30 ；  Y_tikz = 533.77 − y_page
% 线宽：图形本体 0.398 pt（清单 §1.2）。
% 箭头：TikZ arrows.meta 标准箭矢（> = Triangle），箭尖落于清单所给尖端坐标；
%       源图线段本体短于尖端（留白由填充三角补齐），TikZ 由 tip 自动收尾，等效。
% 标签：绝对字号 10.5pt / 行距 12.6pt，inner sep=0（全库既定，勿改相对字号）。
% 要素：7 框 + 6 主干 + 1 回折 = 14 描边 path；7 填充箭头（清单 §1.5 合计 21）。
% ============================================================
\usetikzlibrary{arrows.meta}
\begin{tikzpicture}[
  x=1pt,y=1pt,                             % 画布单位 = 1 pt（本文件坐标一律为 pt）
  every node/.style={font=\fontsize{10.5pt}{12.6pt}\selectfont,inner sep=0pt},
  >={Triangle[length=3.60pt,width=3.36pt]},
  fr/.style={line width=0.398pt},          % 框体
  ln/.style={line width=0.398pt},          % 主干 / 回折线
]

  % ---------- 主干线（x = 216.22 → 36.92；起笔点在段首，箭尖在段尾）----------
  \draw[ln,->] (36.92,184.78) -- (36.92,170.61);   % a 开始 → S=20,i=1
  \draw[ln,->] (36.92,154.01) -- (36.92,139.84);   % b S=20,i=1 → i=2i
  \draw[ln,->] (36.92,125.52) -- (36.92,111.34);   % c i=2i → S=S−i
  \draw[ln,->] (36.92, 95.89) -- (36.92, 81.72);   % d S=S−i → 菱形
  \draw[ln,->] (36.92, 62.82) -- (36.92, 48.64);   % e 菱形下顶点 → 输出 S
  \draw[ln,->] (36.92, 31.38) -- (36.92, 17.21);   % f 输出 S → 结束

  % ---------- 「否」支回折（3 段：右 → 上 → 左），自菱形右顶点右侧 0.28 起笔 ----------
  \draw[ln,->] (74.13,72.27) -- (88.31,72.27) -- (88.31,146.93) -- (36.92,146.93);

  % ---------- 节点框 ----------
  % 1 开始（圆角矩形，r≈3.99）
  \draw[fr,rounded corners=3.99pt] (22.92,184.98) rectangle (50.92,201.82);
  \node at (36.90,193.97) {开始};
  % 2 S=20, i=1
  \draw[fr] (3.01,154.21) rectangle (70.84,170.41);
  \node at (36.925,162.31) {$S=20,\,i=1$};
  % 3 i=2i
  \draw[fr] (20.18,125.72) rectangle (53.67,139.64);
  \node at (36.925,132.68) {$i=2i$};
  % 4 S=S−i
  \draw[fr] (11.15,96.09) rectangle (62.69,111.14);
  \node at (36.92,103.615) {$S=S-i$};
  % 5 i>5?（菱形，右/上/左/下顶点）
  \draw[fr] (73.85,72.27) -- (36.92,81.44) -- (0.00,72.27) -- (36.92,63.10) -- cycle;
  \node at (36.92,72.27) {$i>5?$};
  % 6 输出 S（非直角平行四边形，顶部右斜 6.14）
  \draw[fr] (11.51,31.58) -- (17.65,48.44) -- (62.33,48.44) -- (56.19,31.58) -- cycle;
  \node at (36.92,40.01) {输出 $S$};
  % 7 结束（圆角矩形，r≈3.99）
  \draw[fr,rounded corners=3.99pt] (22.92,0.00) rectangle (50.92,17.01);
  \node at (36.90,9.07) {结束};

  % ---------- 支标（清单 §1.4 实测字形中心）----------
  \node at (45.90,56.37) {是};
  \node at (74.10,64.57) {否};

\end{tikzpicture}
